% Section 3. Drafted from docs/paper/section-3-skeleton.md. Every named result is a gated,
% axiom-free theorem in coq/*.v; statements are read from the source.

\section{The main result}\label{sec:main}

Replicas apply the same events in different orders, because networks delay, reorder and retry
messages, and every replica runs the same repair after each event to restore its invariants. The
question is when replicas that have seen the same events end in the same state. The answer has one
shape at every scale: raw states may differ along the way, and what has to agree is the result after
repair. This section states an exact condition at three scales, a single event
(\cref{sec:main-state}), finite histories (\cref{sec:main-history}) and infinite runs
(\cref{sec:main-recurrence}), together with the single-registry theorem that combines the first two
(\cref{sec:main-single}). The three conditions are not one equation. They form a proved hierarchy,
with counterexamples showing where each level stops (\cref{sec:main-hierarchy}).

Throughout, $X$ is a set of states and $N : X \to X$ is the repair. An event is a map
$e : X \to X$ on raw states; a replica applies an event and then repairs, so its step is
$x \mapsto N(e(x))$.

% ---------------------------------------------------------------------------------------------
\subsection{One event: state descent}\label{sec:main-state}

\begin{definition}[State descent; \coq{DescendsOn}]\label{def:descent}
An event $e$ \emph{descends} on a set $P \subseteq X$ when
\[
  N(e(N(x))) \;=\; N(e(x)) \qquad\text{for every } x \in P.
\]
\end{definition}

The condition says that applying the event to a state or to its repaired form gives the same result
after repair. Repairing early loses nothing the event can observe. For an idempotent repair the definition is the
standard compatibility of a map with the kernel of $N$ (\coq{descends_iff_respects}); what matters here is that
convergence forces it, at the states a system can reach.

\begin{example}[A shelf]\label{ex:shelf}
A shelf holds at most 5 items and has 4 in stock. The events are ``add 3'' and ``remove 1'', and the
repair caps the stock at 5, so $N(x) = \min(x, 5)$. Replica~A applies ``add 3'' (7), repairs (5), then
``remove 1'': it ends at 4. Replica~B applies ``remove 1'' (3), then ``add 3'' (6), and repairs: it
ends at 5. The two replicas saw the same events and ran the same repair, and they disagree.

State descent locates the failure. Take the raw state 7, just after ``add 3'', and let $e$ be
``remove 1''. Then $N(e(N(7))) = N(e(5)) = N(4) = 4$, while $N(e(7)) = N(6) = 5$. The event does not
descend at 7: removing an item before or after the cap is applied gives different answers.
\end{example}

In the single-registry setting of \cref{sec:main-single}, state descent is exactly one of the two
conditions convergence requires. There, the repair is built from a repair step $\rho$ that strictly
decreases a potential $\Phi : X \to \mathbb{N}$ on invalid states, and $N = \rho^{*}$ applies $\rho$
until the state is valid.

\begin{proposition}[\coq{state_descent_iff_cc2}]\label{prop:descent-cc2}
If $\Phi(\rho(s)) < \Phi(s)$ for every invalid $s$, then every event descends on all of $X$ under
$N = \rho^{*}$ if and only if
\[
  \rho^{*}(e(s)) \;=\; \rho^{*}(e(\rho(s))) \qquad\text{for every event } e \text{ and every invalid } s.
\]
\end{proposition}

The right-hand side is the condition called CC2 below: an event does not care whether one repair
step ran before it.

% ---------------------------------------------------------------------------------------------
\subsection{Finite histories: history descent}\label{sec:main-history}

A delivery model says which reorderings of a history count as the same delivery. Under free delivery
every permutation is allowed; under causal delivery, only orders consistent with a happens-before
relation; under at-least-once delivery, also duplicates; under declared independence, swaps of
declared pairs.

\begin{definition}[History descent]\label{def:history}
For a delivery model, given as an equivalence $\equiv$ on histories, and a start $s_0$, the system
satisfies \emph{history descent} when every two equivalent histories give the same result after
repair: $N(\mathrm{run}(h, s_0)) = N(\mathrm{run}(k, s_0))$ whenever $h \equiv k$, where
$\mathrm{run}$ applies the raw events in order.
\end{definition}

The raw runs in this definition never repair. A replica, however, repairs after every event. The next
lemma says that, given state descent, the two views agree, so a checker may work with repaired states
only.

\begin{lemma}[\coq{history_from_descent}]\label{lem:transfer}
Suppose every event descends on $X$. Write $\mathrm{grun}(h, x)$ for the run that applies the events
of $h$ to $x$ and repairs after each one. Then for every equivalence $\equiv$ on histories and every
start $s_0$,
\[
  \bigl(\forall h \equiv k.\; N(\mathrm{run}(h, s_0)) = N(\mathrm{run}(k, s_0))\bigr)
  \iff
  \bigl(\forall h \equiv k.\; \mathrm{grun}(h, N(s_0)) = \mathrm{grun}(k, N(s_0))\bigr).
\]
\end{lemma}

The proof is the observation that, under state descent, $N(\mathrm{run}(w, s_0)) =
\mathrm{grun}(w, N(s_0))$ for every word $w$ (\coq{normalization_descent}). Which reorderings preserve
the result depends on the delivery model; Section~5 gives one theorem pair that decides it for
each model above. Here only free delivery is needed, through the single-registry theorem.

% ---------------------------------------------------------------------------------------------
\subsection{The single-registry theorem}\label{sec:main-single}

A \emph{registry} fixes the states, a validity predicate $\valid$, the events, and the repair step
$\rho$ with its potential $\Phi$. Write $\gov(e) = \rho^{*} \circ e$ for the governed step. A
configuration is a pair $(\sigma, B)$ of a state and a buffer of pending events; a step either applies
a pending event to $\sigma$ or, when $\sigma$ is invalid, applies $\rho$. A normal form is a valid state
with an empty buffer. A state is \emph{reachable} from $s_0$, written $\reach(s_0, \sigma)$, when some
sequence of event and repair steps leads from $s_0$ to $\sigma$.

\begin{theorem}[Compensation commutativity; \coq{cc_exact}]\label{thm:cc}
Let $\Phi(\rho(s)) < \Phi(s)$ for every invalid $s$. For every start $s_0$, every configuration
$(\sigma, B)$ with $\reach(s_0, \sigma)$ has a unique normal form if and only if both of the following
hold at every state $\sigma$ reachable from $s_0$:
\begin{enumerate}[label=(CC\arabic*)]
  \item $\gov(e_2)(\gov(e_1)(\sigma)) = \gov(e_1)(\gov(e_2)(\sigma))$ for all events $e_1, e_2$;
  \item $\gov(e)(\sigma) = \gov(e)(\rho(\sigma))$ for every event $e$, when $\sigma$ is invalid.
\end{enumerate}
\end{theorem}

In words: replicas converge exactly when repaired events commute (CC1) and an event does not care
whether repair ran first (CC2), at the states they can reach. Termination of repair is the standing
hypothesis; it is the third condition of the theory, called WFC in the earlier papers.

The theorem reads directly in the vocabulary of the two descents. CC1 is history descent for free
delivery, stated on the governed steps, and CC2 is state descent on the reachable states.

\begin{corollary}[\coq{nd_derived_from}]\label{cor:descent-form}
Under the hypotheses of \cref{thm:cc}, every configuration $(s_0, B)$ has a unique normal form if and
only if the governed runs from $N(s_0)$ agree on every permutation of every history, and every event
descends on the states reachable from $s_0$.
\end{corollary}

Section~9 shows that this form is derived from the framework's kernels without using
\cref{thm:cc}.

\paragraph{Proof idea, and what is not new.}
Sufficiency is Newman's lemma restricted to the reachable states, with a case split on the two kinds
of critical peak: two events, and an event against a repair step. Necessity uses two small runs. The
buffer $[e_1, e_2]$ in both orders forces CC1; the buffer $[e]$ with the event first versus the repair
first forces CC2. Both ends of each run are normal forms, because $\rho^{*}$ always reaches a valid
state. The result is therefore the critical-pair theorem of Knuth and Bendix and of
Huet~\cite{huet1980}, specialized to a system whose critical pairs are already normal forms: that is
why ``joinable'' becomes ``equal'' and the converse costs nothing. The contribution is the condition
and its necessity at the reachable states, not the proof technique. Two counterexamples fix the
qualifiers: without the restriction to reachable states the converse fails
(\coq{rho_star_qualifier}, \coq{masked_cc1}).

% ---------------------------------------------------------------------------------------------
\subsection{Infinite runs: canonical recurrence}\label{sec:main-recurrence}

When repair may never stop, for example on networks of registries that feed each other around a
cycle, the question becomes one of long-run behavior. A system now has a finite set of actions $J$,
each with an update $u_j : X \to X$, and a schedule $\sigma : \mathbb{N} \to J$ picks the action at
each step. A schedule is \emph{fair} when every action is picked infinitely often, an action that
changes nothing at the current state being a no-op. At this scale $N$ may be any observation $N : X \to C$ of the state, for example its repaired
result. A run \emph{settles} on $c \in C$ when $N(x_n) = c$ for all large enough $n$.

A \emph{bad fair lasso} from $s_0$ for $c$ is a path from $s_0$ to a state $x$ followed by a
nonempty loop from $x$ back to $x$ that picks every action, with some state on the loop outside the
fiber $N^{-1}(c)$.

\begin{theorem}[Canonical recurrence; \coq{C1_closed}]\label{thm:c1}
Let $X$ be finite, and suppose the fiber $N^{-1}(c)$ is closed under every update. Then every fair run
from $s_0$ settles on $c$ if and only if there is no bad fair lasso from $s_0$ for $c$.
\end{theorem}

Settling fails exactly through a fair loop that leaves the target result. The kernel is the classical
fair-cycle characterization of Emerson and Lei~\cite{emersonlei1987}, read through $N$; its role here
is to classify how runs fail (oscillation, livelock, settling on the wrong state) in one vocabulary
(Section~6). Finiteness is needed for the exact form only: a bad fair lasso refutes settling
on any state space (\coq{C1_refutes}).

At this scale the condition reads repair only through its result.

\begin{proposition}[\coq{canon_settles_quotient_invariant}]\label{prop:quotient}
Suppose the updates $u_j$ respect $N$, that is $N(x) = N(y)$ implies $N(u_j(x)) = N(u_j(y))$, and
let $u'_j$ satisfy $N(u'_j(x)) = N(u_j(x))$ for every action and state. If $N(s'_0) = N(s_0)$, then
every fair run of $u$ from $s_0$ settles on $c$ if and only if every fair run of $u'$ from $s'_0$
does.
\end{proposition}

% ---------------------------------------------------------------------------------------------
\subsection{How the conditions relate}\label{sec:main-hierarchy}

The three conditions are not one equation. State descent and history descent are independent of each
other; together, with terminating repair, they are exactly convergence of a single registry
(\cref{cor:descent-form}); and terminating repair connects them to settlement through Newman's lemma
(Section~8). Each separation is witnessed by a mechanized counterexample.

\begin{proposition}\label{prop:separations}
\leavevmode
\begin{enumerate}[label=(\alph*)]
  \item \emph{State descent without history descent} (\coq{descent_not_history}). With $N$ the
    identity every event descends, and two events that do not commute give different results for the
    two orders of one history.
  \item \emph{History descent without state descent} (\coq{history_not_descent}). With one event,
    the successor map on $\mathbb{N}$, and the repair $N$ that sends $0$ and $1$ to $0$ and fixes every
    other number, every reordering of a history is the history itself, so history descent holds; but
    $N(e(N(1))) = N(1) = 0$ while $N(e(1)) = N(2) = 2$.
  \item \emph{Infinite runs need descent} (\coq{trajectory_needs_descent}). There are two systems whose
    updates have the same image under $N$ at every state, neither of which respects $N$, such that,
    from the same start, every fair run of one settles on \textsf{true} and every fair run of the
    other settles on \textsf{false}.
\end{enumerate}
\end{proposition}

\begin{figure}[t]
\centering
\begin{tikzpicture}[
  box/.style={draw, rounded corners, minimum width=3.3cm, minimum height=1.1cm, align=center},
  lab/.style={font=\small, align=center}]
  \node[box] (s) {\textsc{one event}\\ state descent};
  \node[box, right=2.6cm of s] (h) {\textsc{sequences}\\ history descent};
  \node[box, right=2.6cm of h] (c) {\textsc{infinite runs}\\ canonical recurrence};
  \draw[{Latex}-{Latex}, dashed] (s) -- node[lab, above] {neither implies\\ the other} (h);
  \draw[-{Latex}] (h) -- node[lab, above] {with terminating\\ repair} (c);
  \draw ([yshift=-0.35cm]s.south west) -- ([yshift=-0.35cm]h.south east)
    node[lab, midway, below=2pt] {together, with terminating repair: exactly the single-registry theorem};
\end{tikzpicture}
\caption{The three conditions and the proved relations between them
(\cref{prop:separations,cor:descent-form,thm:c1}).}
\label{fig:hierarchy}
\end{figure}

Every condition in this section is exact: necessary and sufficient, with each hypothesis shown to be
needed by a counterexample. Section~4 reads the single-registry theorem against CRDTs, the case
where repair never runs, and Sections~5 to~8 give the
three-layer framework that organizes these conditions across the regimes of the map.


\begin{thebibliography}{9}
\bibitem{huet1980} G.~Huet. Confluent reductions: abstract properties and applications to term
  rewriting systems. \emph{Journal of the ACM} 27(4):797--821, 1980.
\bibitem{emersonlei1987} E.~A.~Emerson and C.-L.~Lei. Modalities for model checking: branching time
  logic strikes back. \emph{Science of Computer Programming} 8(3):275--306, 1987.
\end{thebibliography}
